\begin{proposition}[Proved: derived prime fibres and $K_0$ do not determine the answer]
\label{C:derived_fibre_twins}
Put $D=\mathbb Z_{(2)}$ and
\[
 A=D[v]/(v^2-2),\qquad B=D[u]/(u^2-8).
\]
These are noetherian local domains, finite free of rank two over $D$.
For every nonzero integer $a$ they have
$K_0(A[1/a])=K_0(B[1/a])=\mathbb Z$ with the free module as generator.
Their rational fibres are isomorphic to $\mathbb Q(\sqrt2)$, and their
derived fibres over every prime of $\operatorname{Spec}\mathbb Z$ are
isomorphic as algebras over the corresponding residue field. Nevertheless
$\vartheta_A$ is an isomorphism in every degree, while $\vartheta_B$
fails to be surjective in degree two. More precisely, for
\[
 T=\begin{pmatrix}1+u&2\\-u-4&-1-u\end{pmatrix},\qquad
 M=(B^2,T),
\]
the class $\eta=[M]-[\mathbf1]-[\mathrm{sgn}]$ has order exactly two
and is zero on the rational fibre and at both residue-field points of
$\operatorname{Spec}B$.
\end{proposition}
\begin{proof}
Both displayed monic polynomials are irreducible over $\mathbb Q$:
$2$ and $8$ are not rational squares. Their quotient rings embed in
$\mathbb Q(\sqrt2)$, by $v\mapsto\sqrt2$ and $u\mapsto2\sqrt2$.
The monic presentations also give the stated free $D$-bases.
Since each ring is finite integral over $D$, all its maximal ideals
lie above $(2)$. Both reductions modulo $2$ are
$\mathbb F_2[t]/(t^2)$, which has just one maximal ideal. Hence
$A$ and $B$ are local, with maximal ideals $(v)$ and $(2,u)$
respectively. In the first ring $2=v^2$, so its maximal ideal is
principal. It is a one-dimensional noetherian local domain, because
it is finite integral over the one-dimensional DVR $D$; the usual
principal-maximal-ideal characterization makes $A$ a DVR. Thus
Proposition~\ref{C:DVR} proves the assertion about $\vartheta_A$.

If $a$ is odd, localization at $a$ changes neither ring. If $a$ is
even, it gives their common fraction field, the identification being
$u=2v$. These rings are all local, so their finite projectives are
free and the asserted $K_0$ identifications follow. Both $A$ and
$B$ are flat over $\mathbb Z$: they are free over the flat
$\mathbb Z$-algebra $D$. Consequently all their derived arithmetic
fibres are concentrated in degree zero. At $(0)$ they are the
isomorphic rational fields just specified; at $(2)$ they are both
$\mathbb F_2[t]/(t^2)$; at every odd prime they are the zero algebra,
since that prime is already invertible in $D$. These identifications
also persist after any extension of the residue fields.

In $B$ the identity $u^2=8$ gives
\[
 (1+u)^2+2(-u-4)=1,
\]
so $T^2=I_2$. The ring map
$B\to\mathbb Z/4$, $u\mapsto2$, sends $T$ to
$\left(\begin{smallmatrix}3&2\\2&1\end{smallmatrix}\right)$.
For the invariant of Lemma~\ref{C:Bockstein_parity} this has
$\overline{T-I}=0$ and
\[
 N=\frac{T-I}{2}\bmod2
   =\begin{pmatrix}1&1\\1&0\end{pmatrix},\qquad
 \det(tI-N)=t^2+t+1.
\]
The latter polynomial is irreducible and is fixed by $t\mapsto t+1$.
The corresponding primary module has length one. Its parity is
therefore one on $M$ and zero on each of $\mathbf1$ and
$\mathrm{sgn}$. Scalar extension preserves the admissible exact
sequences, which split on their underlying $B$-modules. It follows
that $\eta\ne0$, and that $[M]$ is outside the degree-two image of
$\vartheta_B$, generated by those two line classes.

The matrix
\[
 J=\begin{pmatrix}1&0\\-1-u&-1\end{pmatrix}
\]
satisfies $J^2=I_2$ and $JT=-TJ$, as direct multiplication shows.
Thus $M\otimes\mathrm{sgn}\simeq M$. The evaluation exact sequence
\textup{(C.1.7)} gives $2[M]=2[\mathbf1]+2[\mathrm{sgn}]$;
together with the parity detector this proves exact order two.
Over the fraction field the trace and determinant of $T$ are $0$
and $-1$, so its two eigenspaces have dimensions one and $\eta$
vanishes there. At the closed point $u=2=0$, the matrix $T$ is the
identity and the trivial and sign lines agree. This proves vanishing
at the remaining residue-field point.
\end{proof}

Thus even the complete family of derived prime fibres, the rational
fibre, and all integer-localized $K_0$ groups fails to distinguish a
positive from a negative answer to C.1. The total rings retain
information absent from these data: for instance the map
$B\to\mathbb Z/4$ used above has no counterpart for $A$, since no
element of $\mathbb Z/4$ has square $2$. This rules out that
particular fibre-only criterion; it does not rule out a condition
involving the structure of the total ring or its higher arithmetic
thickenings.

The finite certificate \texttt{code/C\_finish\_fibre\_twins.py}
checks the integral polynomial matrix identities, the complete
modulo-four quotient map, the common modulo-two fibre, and the
Bockstein operator. Its recorded finite counts and runtime are in
\texttt{logs/C\_finish\_fibre\_twins.json}.
